\section{三类市场的策略适用性差异}

A 股、境内商品/金融期货与虚拟资产市场在交易时间、结算、杠杆、价格限制、交易所结构、账户准入、数据质量和监管边界上均不同。将一套全球机构策略直接横跨三者，会把市场制度风险误写成“分散化”。本章仅描述研究可映射性；它不为任何市场、产品、交易所、账户或跨境路径提供实施许可。

\begin{exhibitbox}[Exhibit 4：全球策略到三类市场的适用性矩阵]
\scriptsize
\begin{tabularx}{\exhibitboxwidth}{L{2.45cm} X X X}
\toprule
\textbf{策略经济引擎} & \textbf{A 股 / 境内公开证券} & \textbf{境内期货} & \textbf{虚拟资产 / 加密市场} \\
\midrule
风险来源与战略配置 & 仅能做公开产品和现金流的风险地图；产品级条款、估值与流动性仍待核验。 & 可把保证金、品种组与风险因子列入共同 NAV；不能把名义金额当作风险。 & 仅作风险研究；波动、托管、稳定币、交易所和 24/7 清算会改变所有结论。 \\
系统化趋势 / CTA & 单股趋势容易与公司/行业风险混合，不能等同于多市场 CTA。 & 在实际合约、结算价、换月、保证金、限仓和交易状态可得时，可做低频纸面研究。 & 经济上可研究趋势，但不构成内地实施路径；数据、场所、托管与合规须逐项验证。 \\
因子 / 风格溢价 & 可做长多质量、价值、动量、投资等公开信息排序；须承认 T+1、涨跌停和无空头腿。 & 可研究跨品种风险分组，但不将期限结构或 carry 误称为无风险收益。 & 资产存活偏差、市场碎片化、资产质量和交易所风险使因子定义更脆弱。 \\
公告 / 事件重定价 & 可研究正式公告后的可得时间、基本面惊喜和流动性；停牌、限价及 T+1 是硬约束。 & 可研究交易所公告、报告和交割相关事件；须区分消息与可交易合约事实。 & 仅研究信息来源、时间戳与操纵风险；不把社媒、KOL 或链上叙事当作可验证事件。 \\
相对价值 / carry & 无稳定借券和对冲时，不称为市场中性或无风险价差。 & 仅在真实双腿、交割、持仓、保证金和成本可以审计时做条件研究。 & 资金费率、基差、借币与场所对手方均可形成尾部风险，不能视作自动收益。 \\
统计套利 / 做市 & 不可等价复制；不存在“相关性高”就可配对交易的推论。 & 不可等价复制；数据/执行/容量不足时不生成信号。 & 不可等价复制；24/7 与碎片化并不会降低基础设施门槛。 \\
\bottomrule
\end{tabularx}
\sourcenote{S01--S15；适用性判断为本报告分析，具体产品/合约/账户仍须以当前规则、合同和权限重新核验。}
\end{exhibitbox}

\section{A 股：长多公开信息研究的现实边界}

A 股研究可以利用上市公司正式披露、交易所公告和公开市场数据，但它应从一开始承认制度差异。普通股票的买入后处置、停复牌、涨跌幅、风险警示、再融资、信息披露节奏、行业结构和流动性都会改变一个看似简单的“入场/止损”规则。对于长多研究，进入之前应将单名、行业和共同因子风险压低到即使发生处置受限也不会破坏共同 NAV。加快卖出规则的书写不能替代这一步。

因此，A 股上的因子和事件框架仅用于候选筛选与公开资料承保。个人研究缺少机构多空、指数期货对冲和借券能力；一日成交额、榜单排名和单一技术指标也不足以证明流动性或执行能力。

\section{境内期货：合约、保证金与交割是策略本身}

期货中的趋势、跨期或相对价值研究必须由实际合约驱动。连续主力图可用于初步观察，却不能代替逐日保存的合约代码、结算价、成交量、持仓量、到期日、合约乘数、交易状态、保证金、涨跌停、限仓与换月记录。字段缺失时，任何回放都无法说明历史某日究竟交易了什么、需要多少保证金、是否可以继续持有、是否受到交割月或持仓限制影响。

本报告只把境内期货定位为\textbf{条件性纸面研究场景}。合约资格、交易所规则、期货公司要求和当前风险承受能力不由本报告判断；若其中任一事实为空，CTA 或价差研究仍为 \qzero{}。

\section{加密市场：研究对象与内地个人实施边界}

加密价格序列可以用于研究趋势、事件扩散、市场微观结构或跨场所风险，但这不自动建立任何投资或交易路径。其独有问题包括全天候交易、市场碎片化、交易所和托管对手方、稳定币、清算瀑布、资金费率、链上/链下数据差异、资产存活偏差、欺诈与操纵风险。即使数据看起来完整，数据来源是否有权使用、事件在何时对不同市场可得、以及一个假设中的订单会落在哪个场所，仍会改变研究结论。

对中国内地居民，本备忘录不把虚拟资产、境外平台、VPN、代持、稳定币转换或任何规避安排视为研究成果或实施方案。与虚拟货币相关的法律、监管、跨境、税务、身份和平台问题应由当期、独立且合格的专业人士判断；本报告不提供法律意见。基于本报告的所有加密相关条目均为 \researchonly{} 与 \qzero{}。

\begin{riskbox}[同一“趋势”不等于同一种风险]
一个趋势信号可以同时出现在股票、期货或代币图表上，但不意味着其风险、成交、结算、保证金、托管与监管可以互换。研究者必须先回答“我到底研究的是哪个市场、哪一份数据、哪一条规则和哪一类结算风险”，再讨论价格行为；不能让相似图形掩盖不同的制度事实。
\end{riskbox}
